\section{The phase-weighted Ces\`aro formula for the peripheral projection}

\begin{lemma}[Orthogonality of unit-modulus phases]
    \label{lem:unit_modulus_phase_cesaro}
    \lean{WeightedCesaro.tendsto_cesaro_phase_sum,
        WeightedCesaro.tendsto_cesaro_geom_zero,
        WeightedCesaro.tendsto_cesaro_geom_one,
        WeightedCesaro.sum_Icc_one_eq_sum_range}
    \leanok
    Let $s$ be a finite set of complex numbers of modulus one and let
    $\lambda\in s$. Then
    \begin{align}
        \lim_{N\to\infty}\frac{1}{N}\sum_{n=1}^{N}\sum_{\mu\in s}
        (\bar\mu\lambda)^{n}
         & =1.
        \label{eq:unit_modulus_phase_cesaro}
    \end{align}
\end{lemma}

\begin{proof}\leanok
    Every ratio $\bar\mu\lambda$ again has modulus one, and
    $\bar\mu\lambda=1$ forces $\lambda=\mu$, since $\mu\bar\mu=1$. The term
    $\mu=\lambda$ therefore equals $1$ for every $n$ and contributes $1$ to
    the mean. For $\mu\neq\lambda$ set $\zeta=\bar\mu\lambda\neq1$; the
    geometric sum
    \begin{align}
        \sum_{n=1}^{N}\zeta^{n}
         & =\frac{\zeta-\zeta^{N+1}}{1-\zeta}
    \end{align}
    is bounded by $2/\lvert1-\zeta\rvert$ uniformly in $N$, so dividing by
    $N$ sends it to zero.
\end{proof}

\begin{lemma}[Ces\`aro means of a vanishing sequence]
    \label{lem:cesaro_mean_of_null_sequence}
    \lean{WeightedCesaro.tendsto_cesaro_zero}
    \leanok
    Let $u_{1},u_{2},\dots$ be a sequence in a complex normed space with
    $u_{n}\to0$. Then
    \begin{align}
        \frac{1}{N}\sum_{n=1}^{N}u_{n}
         & \longrightarrow0.
    \end{align}
\end{lemma}

\begin{proof}\leanok
    The Ces\`aro mean of a convergent sequence converges to the same limit.
\end{proof}

\begin{definition}[Phase-weighted Ces\`aro mean]
    \label{def:weighted_cesaro_mean}
    \lean{Module.End.weightedCesaroMean,
        Module.End.weightedCesaroMean_apply}
    \leanok
    Let $T$ be an endomorphism of a complex vector space and let $s$ be a
    finite set of complex numbers. The \emph{phase-weighted Ces\`aro mean} of
    $T$ over $s$ is
    \begin{align}
        C_{N}(T,s)
         & :=\frac{1}{N}\sum_{n=1}^{N}\sum_{\mu\in s}(\bar\mu T)^{n}.
        \label{eq:weighted_cesaro_mean}
    \end{align}
    Taking for $s$ the eigenvalues of $T$ of modulus one gives the averages of
    \cite[Equation~(6.15)]{Wolf2012Quantum}.
\end{definition}

\begin{lemma}[Action of the phase-weighted means on the spectral splitting]
    \label{lem:weighted_cesaro_mean_on_splitting}
    \lean{Module.End.tendsto_weightedCesaroMean_apply_self_of_mem_iSup_eigenspace,
        Module.End.tendsto_weightedCesaroMean_apply_zero_of_tendsto_pow_zero}
    \leanok
    \uses{def:weighted_cesaro_mean}
    Let $T$ be an endomorphism of a complex normed space and let $s$ be a
    finite set of complex numbers of modulus one.
    \begin{enumerate}
        \item[(i)] On the span of the eigenspaces of $T$ belonging to the
              elements of $s$, the means $C_{N}(T,s)$ converge to the
              identity.
        \item[(ii)] At a vector $X$ with $T^{n}(X)\to0$, the means
              $C_{N}(T,s)(X)$ converge to zero.
    \end{enumerate}
\end{lemma}

\begin{proof}\leanok
    \uses{lem:unit_modulus_phase_cesaro, lem:cesaro_mean_of_null_sequence}
    For (i) it suffices to treat a $\lambda$-eigenvector $X$ with
    $\lambda\in s$, because the span is the sum of the eigenspaces. There
    $(\bar\mu T)^{n}(X)=(\bar\mu\lambda)^{n}X$, so
    $C_{N}(T,s)(X)$ is the scalar
    $\frac{1}{N}\sum_{n=1}^{N}\sum_{\mu\in s}(\bar\mu\lambda)^{n}$ times $X$,
    which tends to $X$ by (\ref{eq:unit_modulus_phase_cesaro}).
    For (ii), multiplication by a unit-modulus scalar does not change norms,
    so $\lVert(\bar\mu T)^{n}(X)\rVert=\lVert T^{n}(X)\rVert\to0$ for every
    $\mu\in s$; the inner sum is finite, hence also tends to zero, and its
    Ces\`aro means tend to zero.
\end{proof}

\begin{theorem}[Phase-weighted Ces\`aro formula for $T_\varphi$]
    \label{thm:peripheral_projection_weighted_cesaro}
    \lean{IsPositiveMap.tendsto_weightedCesaroMean_peripheralProjection,
        IsPositiveMap.tendsto_endEquiv_weightedCesaroMean_peripheralProjection,
        IsPositiveMap.tendsto_weightedCesaroMean_toFinset_peripheralProjection,
        IsPositiveMap.tendsto_endEquiv_weightedCesaroMean_toFinset_peripheralProjection}
    \leanok
    \uses{def:peripheral_spectral_projections, def:weighted_cesaro_mean}
    Let $T:\MN{D}\to\MN{D}$ be positive and trace-preserving, with $D>0$, and
    let $\lambda_{1},\lambda_{2},\dots$ denote its eigenvalues. Then
    \begin{align}
        T_\varphi
         & =\lim_{N\to\infty}\frac{1}{N}\sum_{n=1}^{N}
        \sum_{k:\lvert\lambda_{k}\rvert=1}(\bar\lambda_{k}T)^{n},
        \label{eq:peripheral_projection_weighted_cesaro}
    \end{align}
    both pointwise and in operator norm. This is
    \cite[Equation~(6.15)]{Wolf2012Quantum}.

    \textbf{Local fix (deduplicated phases):} the inner sum runs over the
    distinct peripheral eigenvalues, each phase counted once. Indexed by the
    Jordan blocks of \cite[Equations~(6.4)--(6.5)]{Wolf2012Quantum}, where
    several blocks may share an eigenvalue, the summand
    $(\bar\lambda_{k}T)^{n}$ would repeat each phase according to its
    geometric multiplicity, and for the identity channel the right-hand side
    would converge to $D^{2}$ times the identity rather than $T_\varphi$. The
    comparison is recorded in
    \path{docs/paper-gaps/wolf_ch6_eq615_deduplicated_phases.tex}.
\end{theorem}

\begin{proof}\leanok
    \uses{def:peripheral_spectral_splitting,
        thm:peripheral_spectral_complementarity,
        thm:peripheral_jordan_trivial,
        thm:positive_trace_preserving_mean_ergodic_projection,
        lem:weighted_cesaro_mean_on_splitting,
        lem:tendsto_clm_of_tendsto_apply}
    Split $X=T_\varphi(X)+(X-T_\varphi(X))$ along the complementary spectral
    subspaces. Bounded forward orbits force every eigenvalue to have modulus
    at most one, so every eigenvalue outside the peripheral spectrum has
    modulus strictly below one and $T^{n}$ vanishes on the non-peripheral
    subspace; part (ii) of
    Lemma~\ref{lem:weighted_cesaro_mean_on_splitting} sends the second
    summand to zero. Triviality of the peripheral Jordan blocks identifies
    the peripheral subspace with the span of the peripheral eigenspaces, on
    which part (i) of the same lemma acts as the identity, so the first
    summand converges to $T_\varphi(X)$. Pointwise convergence of
    endomorphisms of a finite-dimensional space is convergence in operator
    norm.
\end{proof}
